% Level 2, Week 11 -- Project Studio 2
\lesson{11}{Project Studio 2}{Tests check the brain of your program, playtesters check the fun, and polish makes it work for everyone.}
\lessoninfo{Concepts: tests for own functions (R), files (R), lists and dictionaries (R), feedback (R), accessibility (R)}{Tools: Thonny, pytest, feedback cards}{CSTA: 2-AP-17, 2-AP-15, 2-IC-21}

\begin{objectives}
  \item The student can write at least 3 pytest tests (normal and edge cases) for the team's own functions and run them. (L2.G6)
  \item The student can add a list, a dictionary and a file that saves and loads to the project (Milestone M3). (L2.G4, L2.G5)
  \item The student can run a playtest, record feedback on cards and turn it into a bug list ordered by importance. (L2.G8)
  \item The student can check the project for accessibility and make at least one improvement (Milestone M4). (L2.G9)
\end{objectives}

\begin{materials}
  \item Print: workbook chapter 11 and homework 11 (one per student); 6 playtest feedback cards per team (cut from workbook chapter 11, or copy them onto index cards).
  \item On each Mac: the team's project folder open in Thonny; pytest already installed (week 8). Check now: \menu{Tools > Manage packages…}, search \texttt{pytest}: it should say \enquote{installed}.
  \item The M3/M4 checklists in the Project Brief. A timer. Prep: 10 minutes reading this plan.
\end{materials}

\needspace{16\baselineskip}
\section{At a glance}
\begin{timeplan}
  \timeslot{0--8}{Warm-up}{M2 round: each team runs its core loop in 30 seconds. Introduce the playtest feedback cards.}
  \timeslot{8--18}{Concept}{Test the brain, not the paint: which functions can pytest check? The main guard, test-only file names, normal and edge cases.}
  \timeslot{18--28}{Guided demo}{You add save/load high score to the skeleton, write \texttt{test\_mygame.py}, run it in the system shell, break the code on purpose, see FAILED, fix it.}
  \timeslot{28--52}{Hands-on build}{28--40: M3 data and 3 tests. 40--46: playtest swap with feedback cards. 46--52: M4 polish, accessibility check, free-build.}
  \timeslot{52--60}{Wrap-up}{Show passing tests, read one kind card aloud, exit ticket, reflection, homework 11.}
\end{timeplan}

\section{Warm-up (0--8 min)}
\begin{unplugged}[Playtest feedback cards]
Each team runs its core loop for 30 seconds. Then hand out the feedback cards and read one aloud.
Each card has four lines: \textbf{I liked…}, \textbf{I was confused when…}, \textbf{A bug I saw:} (what I
did / what happened / what I expected), \textbf{One idea…}. Practise with Star Catcher or a team's
game: everyone fills one card in 2 minutes. Say: \enquote{A good card helps the maker. \enquote{It's bad}
does not help. \enquote{When I pressed space nothing happened, I expected a jump} helps a lot.}
\end{unplugged}

\section{Concept (8--18 min)}
\begin{teacherscript}
“In week 8 you tested functions I gave you. Today you test your own. Pytest cannot see whether a
star looks nice, but it can check a function that takes values in and gives an answer back. So we
test the brain, not the paint: \texttt{next\_x}, \texttt{touching}, \texttt{add\_score}, \texttt{line\_to\_shape}.

Two tricks. One: the start-up code sits under \texttt{if \_\_name\_\_ == "\_\_main\_\_":}, so pytest
can borrow your functions without starting the game. Two: file functions take the file name as a
parameter. The game uses \texttt{highscore.txt}, the tests use \texttt{test\_scores.txt}, so tests never
wipe your real high score. Test one normal case and some edge cases, because bugs hide at the edges:
the edge of the screen, a score of 0, a file that does not exist yet.”
\end{teacherscript}

\begin{analogy}[Car factory checks]
A car factory tests the brakes on a machine: press the pedal, measure the stop. That is a unit test,
like pytest. Then a real driver takes the car round the track and says \enquote{the seat is
uncomfortable}. That is a playtest. You need both: machines check the brain, people check the fun.
\end{analogy}

\section{Guided demo (18--28 min)}
\begin{workedexample}[Save, load and three tests]
\begin{enumerate}
  \item Open last week's \texttt{mygame.py}. Under \texttt{next\_x}, type the two file functions from week 7 (code A below). Say: \enquote{No file yet means a high score of 0.}
  \item \menu{File > New}, save as \texttt{test\_mygame.py} in the same folder (the name must start with \texttt{test\_}) and type code B: one normal case, one edge case, and two file tests that use a test-only file name.
  \item Click in \texttt{mygame.py}, press \keys{F5} once (this tells Thonny which folder you are in), close the game window, then \menu{Tools > Open system shell…}. Type \texttt{ls} and press \keys{Return}: both files should be listed.
  \item Type the command in box C. (If it says \enquote{command not found}, type \texttt{python3} instead of \texttt{python}.) Read the output together: four PASSED lines, \enquote{4 passed}.
  \item Break it on purpose: in \texttt{next\_x} change the last line to \texttt{return new\_x} (no clamp), save, and press \keys{↑} then \keys{Return} in the shell to run the tests again. Say: \enquote{The test caught it. 300 is off the screen; we expected 280. FAILED is our friend.} Undo with \keys{Cmd}+\keys{Z}, save, run again: all PASSED.
\end{enumerate}
\end{workedexample}

\begin{pycode}[A. Added to mygame.py (tested)]
def load_high_score(filename):
    """Read the best score from a file. No file yet? Then it is 0."""
    try:
        with open(filename) as f:
            return int(f.read().strip())
    except FileNotFoundError:
        return 0

def save_high_score(score, filename):
    """Write the best score to a file (replaces the old one)."""
    with open(filename, "w") as f:
        f.write(str(score))
\end{pycode}
\begin{pycode}[B. test\_mygame.py (tested)]
from mygame import load_high_score, next_x, save_high_score

def test_moves_right_normally():
    assert next_x(0, 1) == 20

def test_stops_at_the_right_edge():
    assert next_x(280, 1) == 280

def test_high_score_round_trip():
    save_high_score(42, "test_scores.txt")
    assert load_high_score("test_scores.txt") == 42

def test_no_file_means_zero():
    assert load_high_score("no_such_file.txt") == 0
\end{pycode}
C. In the system shell, type the first line; pytest prints the rest:
\begin{shellcode}
python -m pytest test_mygame.py -v
test_mygame.py::test_moves_right_normally PASSED        [ 25%]
test_mygame.py::test_stops_at_the_right_edge PASSED     [ 50%]
test_mygame.py::test_high_score_round_trip PASSED       [ 75%]
test_mygame.py::test_no_file_means_zero PASSED          [100%]
=================== 4 passed in 0.08s ===================
\end{shellcode}
After breaking \texttt{next\_x} on purpose:
\begin{outputcode}
test_mygame.py::test_stops_at_the_right_edge FAILED     [ 50%]
E       assert 300 == 280
E        +  where 300 = next_x(280, 1)
FAILED test_mygame.py::test_stops_at_the_right_edge - assert 300 == 280
============== 1 failed, 3 passed in 0.07s ==============
\end{outputcode}

\section{Hands-on build (28--52 min)}
\begin{pairroles}[Tester and coder]
\textbf{28--40, M3:} one partner writes tests while the other adds the data (a list, a settings
dictionary, the saved file). Swap at minute 34. \textbf{1 student:} you write test names together
on paper first, then the student codes both. \textbf{3:} coder, tester and \enquote{file keeper} who
checks that \texttt{highscore.txt} (or the saved design) appears in the folder and loads again.
\textbf{4:} two teams of two.
\end{pairroles}

Test ideas: arcade teams test movement at both edges, \texttt{touching} near and far, scoring, and
save/load; Art Studio teams test turning a shape into a line of text and back, and a save-then-load
round trip (example in the answer key, week 11). \textbf{40--46, playtest swap:} teams play each other's
project for 4 minutes and fill in 2 feedback cards each. With 1 or 2 students, you (or another family
member) are the playtester.

\begin{bughunt}[Playtest bug list]
Back at their own computer, each team sorts its cards into a bug list on paper with three columns:
\textbf{What I did}, \textbf{What happened}, \textbf{What I expected}. Mark each line \textbf{MUST}
(crashes or cannot play), \textbf{SHOULD} (confusing) or \textbf{COULD} (nice idea). Fix the MUST lines
first. If a bug is in a logic function, write a test that shows it (it fails), then fix the code
(it passes). That test now guards the bug forever.
\end{bughunt}

\textbf{46--52, M4 polish.} Each team does the accessibility review below and makes at least one change.

\begin{safety}[Accessibility review of our own game]
Read each line aloud; the team ticks it or fixes it:
\begin{checklist}
  \item Colours come from the Okabe--Ito palette (\texttt{\#E69F00}, \texttt{\#56B4E9}, \texttt{\#009E73}, \texttt{\#0072B2}, \texttt{\#D55E00}, \texttt{\#CC79A7}…).
  \item Nothing depends on colour alone: good and bad things also differ in shape or size, or have a label.
  \item Text is big enough: \texttt{font=("Arial", 18, "normal")} or larger.
  \item The keys are explained on the screen (\enquote{← → to move, space to start}).
  \item Nothing flashes quickly (more than 3 times a second can hurt some people's eyes or trigger seizures).
  \item The game is not too fast to start: there is an easy setting or a slow first level.
\end{checklist}
\end{safety}

\begin{freebuild}[One polish touch]
Any time left: add a title, a \enquote{how to play} line, a nicer player shape, a start key, or an
end message that shows the high score. Polish only after the tests pass.
\end{freebuild}

\section{Wrap-up (52--60 min)}
Each team shows its pytest output and reads one kind card aloud. Tick M3 and M4; hand out homework 11.

\begin{formativecheck}[Exit ticket]
\textbf{Look for:} at least 3 tests that call the team's own functions and pass; the program saves
something to a file and reads it back; one accessibility change made.
\textbf{Ask:} “Why does our test use \texttt{test\_scores.txt} instead of \texttt{highscore.txt}?”
\textbf{Answer:} so running the tests never overwrites the real high score.
\end{formativecheck}

\begin{reflection}
Which test are you proudest of, and which bug did it catch (or could it catch)? What was the most
useful playtest card? What did you change so more people can enjoy your project?
\end{reflection}

\section{Notes for the teacher}
\begin{commonmistakes}
  \item \textbf{The game window opens when pytest runs}, or pytest hangs. The start-up code is not under \texttt{if \_\_name\_\_ == "\_\_main\_\_":}. Move it there.
  \item \textbf{\enquote{collected 0 items}.} The file name or function name does not start with \texttt{test\_}, or pytest was run in the wrong folder (type \texttt{ls} to check).
  \item \textbf{\enquote{ModuleNotFoundError: No module named mygame}.} The test file is in a different folder from the game, or the game file has a different name.
  \item \textbf{Tests that cannot fail} (\texttt{assert True}, or no \texttt{assert} at all). Ask: \enquote{Break your code on purpose. Does this test turn red?}
\end{commonmistakes}

\begin{differentiation}
\supportpath Give three test names to fill in: \texttt{test\_moves\_normally}, \texttt{test\_stops\_at\_edge},
\texttt{test\_no\_file\_means\_zero}, all for functions the team already has. Copying the shape of
the example tests is fine.
\extensionpath Write 5 or more tests, including one written \emph{before} the code (test first),
or keep a top-5 high-score list in the file and test that a new score lands in the right place.
\end{differentiation}

\begin{stretchchallenge}[Top five]
Write \texttt{add\_score(scores, new\_score)} that returns a new list of the best five scores,
biggest first. Test it with an empty list, a full list where the new score is too small, and a
new score that should go to the top. Save the list in the file, one score per line.
\end{stretchchallenge}

\begin{teachernote}[If you are short on time]
Shorten the playtest swap to 3 minutes with one card each, and do the accessibility review as a
quick read-aloud in the wrap-up. Do not cut the 3 tests: they are a project requirement. Teams
that are behind can finish M4 polish at home (homework 11) or in the first 5 minutes of week 12.
\end{teachernote}
